\documentclass[11pt,a4paper]{article}

\usepackage[T1]{fontenc}
\usepackage[utf8]{inputenc}
\usepackage{lmodern}
\usepackage{geometry}
\usepackage{amsmath,amssymb,amsfonts}
\usepackage{booktabs}
\usepackage{array}
\usepackage{graphicx}
\usepackage{hyperref}
\usepackage{enumitem}
\geometry{margin=2.2cm}

\title{Fractional Teleoperation in ROS 2:\\Scientific and Algorithmic Description of the Node Implementation}
\author{Extender Workspace -- fractional\_teleoperation\_node}
\date{\today}

\begin{document}
\maketitle

\begin{abstract}
This note documents the mathematical model and the exact implemented algorithm of the ROS 2 standalone node
\texttt{fractional\_teleoperation\_node}. The control law is based on a fractional-order relation between joystick input
and Cartesian offset, discretized with a finite-memory Gr\"unwald--Letnikov (GL) recursion.
The document first introduces the scientific context of fractional analysis, then details dynamic-order adaptation,
ramp-based offset scaling, reference drift (first-order and fractional modes), and snap-on-release logic.
The goal is scientific consistency with the implementation and operational clarity for tuning and validation.
\end{abstract}

\section{Introduction: Why Fractional Analysis for Teleoperation}
Fractional calculus extends differentiation and integration to non-integer order and enables operators with memory,
which is attractive in human-in-the-loop systems where command smoothness and motion persistence are important
\cite{oldham_spanier_1974,podlubny_1999,monje_2010,diethelm_2010}.

In classical integer-order teleoperation mappings, response shaping is mostly obtained from static gains and filters.
A fractional operator introduces history-dependent behavior directly in the law:
\begin{equation}
{}^{C}_{0}D_t^{\alpha} x(t) = f(t), \quad \alpha \in (0,2),
\end{equation}
where $\alpha$ controls memory depth and dynamic character. In practice:
\begin{itemize}[leftmargin=*]
\item $\alpha \approx 1$ behaves close to an ordinary integrator relation.
\item Lower $\alpha$ yields stronger memory and persistent motion trends.
\item Variable $\alpha$ enables context-sensitive response (fine motion at low input, faster motion at high input).
\end{itemize}

The node implementation uses a finite-memory GL recursion for real-time operation, a common pragmatic approximation
in control and signal-processing implementations \cite{podlubny_1999,oustaloup_2000}.

\section{Implemented Control Objective}
Let $u_k \in \mathbb{R}^3$ be the sampled joystick input (linear and angular channels treated separately),
$\Delta x_k$ the fractional offset, and $x_{\mathrm{ref},k}$ a drifting reference position.
The implemented law is:
\begin{equation}
{}^{C}_{0}D_t^{\alpha} \Delta x(t) = K\,u(t),
\end{equation}
with reconstructed desired state
\begin{equation}
x_{\mathrm{des},k} = x_{\mathrm{ref},k} + S_k\,\Delta x_k,
\end{equation}
where $S_k$ is a time-varying ramp scale (separate maxima for linear and angular channels).
The published velocity command is computed from the scaled offset derivative:
\begin{equation}
v_k = \frac{S_k\Delta x_k - S_{k-1}\Delta x_{k-1}}{\Delta t}.
\end{equation}

\section{Discrete Fractional Recursion (GL Approximation)}
For memory length $L$ and sample time $\Delta t$, coefficients are generated recursively as:
\begin{equation}
c_0 = 1, \qquad
c_k = c_{k-1}\frac{k-1-\alpha}{k}, \quad k \ge 1.
\end{equation}

The node update uses:
\begin{equation}
\Delta x_k = (\Delta t)^{\alpha}K u_k - \sum_{j=1}^{L'} c_j \Delta x_{k-j},
\end{equation}
where $L' = \min(L,\text{available history})$. This is applied independently to linear and angular vectors.

\paragraph{Finite-memory implication.}
The GL operator is truncated to $L$ samples; therefore the implementation is an approximation of ideal long-memory
fractional dynamics. Increasing $L$ improves fidelity at increased computational cost.

\section{Dynamic Fractional Order \texorpdfstring{$\alpha$}{alpha}}
When enabled, $\alpha$ is adapted from joystick norm $\lambda = \lVert u_{\text{linear}} \rVert$:
\begin{equation}
\alpha(\lambda)=
\begin{cases}
\alpha_{\min}, & \lambda < l_0,\\
\alpha_{\max}, & \lambda > l_{\max},\\
\alpha_{\min} + (\alpha_{\max}-\alpha_{\min})\dfrac{\lambda-l_0}{l_{\max}-l_0}, & \text{otherwise}.
\end{cases}
\end{equation}

When $|\alpha_{\text{new}}-\alpha_{\text{last}}| > 10^{-3}$, GL coefficients are recomputed.

\paragraph{Implementation caveat (variable order).}
The history buffers are not reprojected when $\alpha$ changes; only coefficients are refreshed.
This is an engineering approximation for real-time operation and should be interpreted accordingly
for strict variable-order fractional analysis.

\section{Gain Normalization Across Fractional Orders}
If normalization is enabled, $K$ is recomputed each cycle using one of two modes:
\begin{align}
\texttt{dt mode:} \quad & K = v_{\max}(\Delta t)^{1-\alpha}, \\
\texttt{perceptual mode:} \quad & K = v_{\max}\,\frac{\Gamma(\alpha)}{t_{\mathrm{ref}}^{\alpha-1}}.
\end{align}
For numerical safety in perceptual mode, a lower bound is used before evaluating $\Gamma(\alpha)$.

\section{Offset Ramp Scaling}
Joystick activity is detected by $\lVert u_{\text{linear}}\rVert > \tau_{\text{active}}$.
An active-duration timer $T_{\text{active}}$ is incremented while active and reset to zero otherwise.
Define
\begin{equation}
r_k =
\begin{cases}
\min\left(1,\dfrac{T_{\text{active}}}{T_{\text{ramp}}}\right), & T_{\text{ramp}} > 0,\\
1, & T_{\text{ramp}} = 0.
\end{cases}
\end{equation}
Then a profile function $\phi(\cdot)$ produces $\phi(r_k)$ and scales become
\begin{equation}
S^{\text{lin}}_k = S^{\text{lin}}_{\max}\phi(r_k), \qquad
S^{\text{ang}}_k = S^{\text{ang}}_{\max}\phi(r_k).
\end{equation}

Supported profiles include linear, quadratic, smoothstep, sigmoid, exponential, inverse exponential,
tanh, cubic Hermite, and step.

\section{Reference Drift and Snap-On-Release}
\subsection{Reference Drift Gate}
Drift is applied only if enabled and
\begin{equation}
\lVert u_{\text{linear}}\rVert \ge \tau_{\text{drift}}.
\end{equation}

\subsection{First-Order Reference Drift}
In first-order mode:
\begin{equation}
x_{\mathrm{ref},k+1} = x_{\mathrm{ref},k} + \beta\,\Delta t\,(x_{\mathrm{des},k}-x_{\mathrm{ref},k}),
\end{equation}
with clamped blend $\beta\Delta t\in[0,1]$ in code.

\subsection{Fractional Reference Drift}
In fractional mode, define error $e_k = x_{\mathrm{des},k}-x_{\mathrm{ref},k}$ and use a second fractional integrator:
\begin{align}
f_k &= \mathcal{F}_{\alpha_r,K_r}(e_k),\\
x_{\mathrm{ref},k+1} &= x_{\mathrm{ref},k} + \Delta t\, f_k,
\end{align}
where $\mathcal{F}_{\alpha_r,K_r}$ is the same GL-based operator class with parameter mapping
\begin{equation}
(\alpha_r,K_r)=\left(\texttt{reference\_fractional\_alpha},\texttt{reference\_fractional\_gain}\right).
\end{equation}

\subsection{Snap-On-Release Transition}
Let $a_k$ be joystick active state ($a_k\in\{0,1\}$). If snap is enabled and transition
$ a_{k-1}=1 \rightarrow a_k=0$ occurs, the node executes:
\begin{align}
x_{\mathrm{ref},k} &\leftarrow x_{\mathrm{ref},k} + S_{k-1}\Delta x_{k^-},\\
\Delta x_k &\leftarrow 0, \\
\text{fractional histories} &\leftarrow \emptyset,
\end{align}
for both linear and angular channels.
This makes joystick release recenter the fractional state around the current desired position.

\section{Algorithmic Form (Node Update Loop)}
At each timer tick:
\begin{enumerate}[leftmargin=*]
\item Read latest joystick linear/angular vectors and operation mode.
\item Update dynamic $\alpha$ (if enabled); refresh GL coefficients if threshold exceeded.
\item Update gain $K$ (if adaptive gain normalization enabled).
\item Save previous scaled offsets for derivative computation.
\item Apply GL fractional integration to linear and angular channels to update $\Delta x_k$.
\item Update ramp timer from activity condition and compute profile-based scales.
\item Form scaled offsets and desired state $x_{\mathrm{des},k}=x_{\mathrm{ref},k}+S_k\Delta x_k$.
\item If snap condition is met (active-to-inactive transition), shift reference by previous scaled offset,
clear fractional offsets and histories.
\item If reference drift gate is open, update reference with first-order or fractional mode.
\item Compute velocities from scaled-offset derivative.
\item Transform to selected frame (base or end-effector), apply translation/rotation mode gating,
then apply output velocity scale.
\item Publish velocity command and optional markers.
\end{enumerate}

\section{Runtime Parameter Set Used in Current Configuration}
The active YAML in this workspace uses the following key values:
\begin{center}
\begin{tabular}{>{\raggedright\arraybackslash}p{0.46\linewidth} >{\raggedright\arraybackslash}p{0.44\linewidth}}
\toprule
Parameter & Value \\
\midrule
$\alpha$ & 0.5 \\
$K$ (manual, no normalization) & 1.0 \\
$L$ (memory length) & 50 \\
$\Delta t$ & 0.01 s \\
Dynamic alpha & enabled, $\alpha_{\min}=0.1$, $\alpha_{\max}=1.0$, $l_0=0.5$, $l_{\max}=1.0$ \\
Ramp scaling & enabled via $S^{\text{lin}}_{\max}=2.0$, $S^{\text{ang}}_{\max}=1.0$, $T_{\text{ramp}}=2.5$ s, profile exponential \\
Snap on release & enabled \\
Reference drift & enabled, threshold $\tau_{\text{drift}}=0.5$, mode fractional, $\alpha_r=0.5$, $K_r=0.5$ \\
\bottomrule
\end{tabular}
\end{center}

\section{Reproducible Analysis Outcomes (Analysis Folder)}
An executable analysis pipeline is provided in the package analysis folder:
\begin{itemize}[leftmargin=*]
\item model and primitives: \texttt{analysis/fractional\_node\_model.py}
\item runner: \texttt{analysis/run\_analysis.py}
\item baseline outputs: \texttt{analysis/results/}
\end{itemize}

The baseline command is:
\begin{equation}
\texttt{/bin/python3 run\_analysis.py}
\end{equation}
executed from \texttt{src/controllers/fractional\_teleoperation/analysis}.

Generated outputs include:
\begin{itemize}[leftmargin=*]
\item linear fixed-$\alpha$ stability table and spectral-radius plot,
\item controllability rank and finite-horizon Gramian table,
\item bounded-input (ISS-style) response table,
\item hybrid threshold sweep table and sensitivity plot,
\item $\alpha$--memory sweep table and heatmap.
\end{itemize}

For the current baseline:
\begin{itemize}[leftmargin=*]
\item Spectral radius increases with $\alpha$ and reaches 1.0 at $\alpha=1.0$ (marginal in this discrete linear check), while remaining below 1 for $\alpha \in [0.1,0.9]$.
\item Controllability rank equals state dimension (49) on the tested fixed-$\alpha$ augmented models, while Gramian conditioning degrades with increasing $\alpha$.
\item In bounded random-input tests ($\|u\|_\infty \le 1$), peak velocity magnitude decreases strongly with $\alpha$ (about 357 at $\alpha=0.1$ to about 3.29 at $\alpha=1.0$), while reference excursion increases.
\item In hybrid threshold sweeps, increasing \texttt{joystick\_active\_threshold} from 0.005 to 0.05 increases observed transition count (99 to 274) and snap events (49 to 137), with mean dwell time decreasing (12.0 to 4.36 steps).
\item In the $\alpha$--memory sweep, peak velocity is primarily $\alpha$-driven (weak dependence on memory length), but reference excursion grows with memory length.
\end{itemize}

These results are implementation-faithful computational checks, not formal global proofs for the full variable-order hybrid system.

\section{Scientific and Engineering Notes}
\begin{itemize}[leftmargin=*]
\item The finite-memory GL update is a practical approximation; sensitivity to memory length should be evaluated experimentally.
\item With variable $\alpha$, coefficient refresh without history re-identification is computationally efficient but not a strict variable-order discretization.
\item Perceptual gain normalization can produce large gains at very low $\alpha$ because of $\Gamma(\alpha)$ behavior; practical lower bounds and safety saturation remain essential.
\item Drift and snap combine to reduce long-term offset accumulation and improve joystick recentering ergonomics.
\end{itemize}

\section{Scope Boundary}
This document targets the standalone node implementation only.
The ROS 2 controller plugin in the same package shares core fractional routines but does not include all node-specific
behavioral blocks (notably the same drift/snap/ramp orchestration).

\begin{thebibliography}{99}

\bibitem{oldham_spanier_1974}
K. B. Oldham and J. Spanier,
\emph{The Fractional Calculus}.
Academic Press, 1974.

\bibitem{podlubny_1999}
I. Podlubny,
\emph{Fractional Differential Equations}.
Academic Press, 1999.

\bibitem{monje_2010}
C. A. Monje, Y. Chen, B. M. Vinagre, D. Xue, and V. Feliu,
\emph{Fractional-Order Systems and Controls: Fundamentals and Applications}.
Springer, 2010.

\bibitem{diethelm_2010}
K. Diethelm,
\emph{The Analysis of Fractional Differential Equations}.
Springer, 2010.

\bibitem{oustaloup_2000}
A. Oustaloup, F. Levron, B. Mathieu, and F. M. Nanot,
``Frequency-band complex noninteger differentiator: characterization and synthesis,''
\emph{IEEE Transactions on Circuits and Systems I: Fundamental Theory and Applications},
vol. 47, no. 1, pp. 25--39, 2000.

\bibitem{valerio_sa_da_costa_2013}
D. Val\'erio and J. S\'a da Costa,
\emph{An Introduction to Fractional Control}.
Institution of Engineering and Technology, 2013.

\end{thebibliography}

\end{document}
